% section: 二項定理

\section{二項定理}\label{節：二項定理}
非負の整数$n$と実数$x,y$について
\[
  (x+y)^n=\sum_{k=0}^n\binom nk x^{n-k}y^k
\]
が成り立つ.これを\bookterm{binomial}{二項定理}という.\footnote{この多項式の式で0乗は因子を含まない積を表し,底が$0$でも$1$と読む.}

$\binom nk={}_n\mathrm C_k$は$n$個から$k$個を選ぶ組合せの数である.$0\le k\le n$では
\BookMathLayout{\[
    \binom nk=\frac{n!}{k!(n-k)!}
    =\frac{n(n-1)\cdots(n-k+1)}{k!}
  \]}{\[
    \begin{aligned}
      \binom nk&=\frac{n!}{k!(n-k)!}\\
      &=\frac{n(n-1)\cdots(n-k+1)}{k!}.
    \end{aligned}
  \]}
$0!=1$とする.$k=0$は何も選ばない一通りなので$\binom n0=1$である.
範囲外の整数$k<0$や$k>n$では$\binom nk=0$と定め,負の整数を階乗表示へ代入しない.
例えば$\binom12=0$であり,$\binom n2=\dfrac{n(n-1)}{2}$は$n=0,1$でも使える.
\booknavref{節：総和型}{総和型}や\booknavref{節：二項係数で多項式を表す}{二項係数と差分}でこの表示を使う.

係数の理由を考える.$(x+y)^n$を展開して$x^{n-k}y^k$を作るには,$n$個の因子のうち$k$個から$y$を,残りから$x$を選ぶ.
その選び方が$\binom nk$通りなので,これを二項係数と呼ぶ.
例えば三つの因子では$x^2y$も$xy^2$も三通りあり,
\[
  (x+y)^3=x^3+3x^2y+3xy^2+y^3.
\]
係数はパスカルの三角形の3段目にも並ぶ.複素数にも同じ定理が成り立つ.
発展演習の数え上げにも使う.
